% GENERATED FILE. Edit canonical lessons, then run npm run generate:books.
% canonical-source-hash: fnv1a64:fdb4cfbcfda80aa0
% canonical-lessons: ZH-C478-tangshang, ZH-C478-shangkou, ZH-W478-shang-hurt, ZH-C478-zhongshang, ZH-C478-niushang, ZH-C478-baoza

\chapter{Injuries}
\label{ch:zh-injuries}
\hlchaptermodality{478}

\begin{chapteropening}
\textbf{By the end of this chapter:} \emph{I can say I burned myself, talk about a wound and serious injuries, a sprain and having a wound dressed.}

Complete the last lesson of chapter 478: I can say I burned myself, talk about a wound and serious injuries, a sprain and having a wound dressed.
\end{chapteropening}

\section[tàngshāng]{tàngshāng — to scald; to burn (with heat)}

\label{lesson:ZH-C478-tangshang}



\begin{quote}
Before the new one: say the Mandarin for one's digestion, then the Mandarin for an appetite.
\end{quote}

\subsection*{You'll want to know: tàngshāng}
\emph{tàngshāng} — \textquotedblleft{}to scald; to burn (with heat)\textquotedblright{}.

\begin{quote}
\emph{wǒ zuò fàn de shíhou tàngshāng le shǒu} — I burned my hand while I was cooking.
\end{quote}

Its characters come later; for now it is a word for the ear.

\begin{cousinweb}
\emph{tàngshāng} is built from \emph{tàng} 'scalding hot' and \emph{shāng} 'to hurt'. The second is inherited from Middle Chinese \emph{syang} (Guangyun, in Baxter's transcription), glossed 'to harm'; no dictionary we use records an older form of the first under this spelling, so its history is not settled.
\end{cousinweb}

\begin{grammarlens}[title={What you've built}]
One more word for saying how it was.
\end{grammarlens}

\subsection*{Your turn}
\begin{itemize}
\raggedright
  \item \emph{Say it:} \emph{tàngshāng}
  \item \emph{Say it:} \emph{tàngshāng}, once more
  \item \emph{Say it:} \emph{wǒ zuò fàn de shíhou tàngshāng le shǒu} — I burned my hand while I was cooking.
  \item \emph{Recall:} say the Mandarin for one's digestion, then the Mandarin for an appetite, then say \emph{tàngshāng} again
\end{itemize}

\begin{tcolorbox}[breakable,colback=teal!4,colframe=teal!35!black,title={Before you move on}]
What does tàngshāng mean? (\textquotedblleft{}To scald; to burn (with heat)\textquotedblright{}.) Say it once more.
\end{tcolorbox}
\section[\texorpdfstring{\zh{伤口} (shāngkǒu)}{伤口 (shāngkǒu)}]{\zh{伤口} (shāngkǒu) — wound; cut}

\label{lesson:ZH-C478-shangkou}



\begin{quote}
Before the new one: say the Mandarin for an appetite, then the Mandarin for to scald oneself.
\end{quote}

\subsection*{You'll want to know: \zh{伤口}}
\textbf{\zh{伤口}} — \emph{shāngkǒu} — \textquotedblleft{}wound; cut\textquotedblright{}.

\begin{quote}
\textbf{\zh{我感觉伤口很疼体温也高了}} \emph{wǒ gǎnjué shāngkǒu hěn téng tǐwēn yě gāo le} — My wound really hurts, and my temperature has gone up too.
\end{quote}

The new character, in pieces: \zh{伤} — six strokes in three pieces: on the left \zh{亻}, two strokes; at the top right a short stroke falling to the left and a bar; and under them \zh{力} \emph{lì}, two strokes. The next lesson puts it in your hand.

\begin{cousinweb}
\textbf{\zh{伤口}} (\emph{shāngkǒu}) is built from \zh{伤} 'to hurt' and \zh{口} 'an opening', both inherited from Middle Chinese, \emph{syang} and \emph{khuwX} (Guangyun, in Baxter's transcription), which glosses the first 'to harm' and the second 'the mouth'.
\end{cousinweb}

\begin{grammarlens}[title={What you've built}]
One more word for saying how it was.
\end{grammarlens}

\subsection*{Your turn}
\begin{itemize}
\raggedright
  \item \emph{Say it:} \emph{shāngkǒu}
  \item \emph{Say it:} \emph{shāngkǒu}, once more
  \item \emph{Say it:} \emph{wǒ gǎnjué shāngkǒu hěn téng tǐwēn yě gāo le} — My wound really hurts, and my temperature has gone up too.
  \item \emph{Recall:} say the Mandarin for an appetite, then the Mandarin for to scald oneself, then say \emph{shāngkǒu} again
\end{itemize}

\begin{tcolorbox}[breakable,colback=teal!4,colframe=teal!35!black,title={Before you move on}]
What does \zh{伤口} mean? (\textquotedblleft{}Wound; cut\textquotedblright{}.) Say it once more.
\end{tcolorbox}
\section[\texorpdfstring{\zh{伤} (shāng)}{伤 (shāng)}]{\zh{伤} — \zh{亻}, then a short top over \zh{力}}

\label{lesson:ZH-W478-shang-hurt}



\begin{quote}
The word first, then the character inside it.

\begin{itemize}
\raggedright
  \item \emph{Recall:} say \emph{shāngkǒu} — \textbf{R1}, one lesson back, before it decays
\end{itemize}

\begin{quote}
\zh{力} \zh{你}
\end{quote}

\emph{Read it:} each one aloud, with nothing said first — \emph{lì}, \emph{nǐ}
\end{quote}

\subsection*{Script — \zh{亻}, then a short top over \zh{力}}
\hlblockfigure{\detokenize{figures/ZH-W478-shang-hurt-filmstrip.pdf}}{How \zh{伤} is written, stroke by stroke}

\begin{quote}
\zh{伤}
\end{quote}

\textbf{Three pieces: the left, then the top right, then the bottom right.} Count the pieces, not the strokes.

\textbf{The left piece, \zh{亻}, two strokes.} A stroke falling to the left, then a vertical down.

\textbf{The top right, two strokes.} A short stroke falling to the left, then a bar from left to right beside it.

\textbf{The bottom right, \zh{力}, two strokes.} A stroke that runs right, turns down and hooks to the left, then a long stroke falling to the left through it.

2 + 2 + 2 = six.

\subsection*{Your turn}
\begin{itemize}
\raggedright
  \item \emph{Write it:} \zh{伤}, all six, with the model in front of you
  \item \emph{Cover:} write it again from memory, and count the pieces, not the strokes
\end{itemize}

\begin{tcolorbox}[breakable,colback=teal!4,colframe=teal!35!black,title={Before you move on}]
How many pieces does \zh{伤} have? (\textbf{Three, two strokes each}.) Which written character ends it? (\textbf{\zh{力}}.)

Source: \href{https://raw.githubusercontent.com/chanind/hanzi-writer-data/68d10a4b21150cae5e1ebbd223eed289cf32d90c/data/\%E4\%BC\%A4.json}{Hanzi Writer Data \zh{伤}, pinned stroke paths and medians}.
\end{tcolorbox}
\section[\texorpdfstring{\zh{重伤} (zhòngshāng)}{重伤 (zhòngshāng)}]{\zh{重伤} (zhòngshāng) — serious injury; badly hurt}

\label{lesson:ZH-C478-zhongshang}



\begin{quote}
Before the new one: say the Mandarin for to scald oneself, then the Mandarin for an open wound.
\end{quote}

\subsection*{You'll want to know: \zh{重伤}}
\textbf{\zh{重伤}} — \emph{zhòngshāng} — \textquotedblleft{}serious injury; badly hurt\textquotedblright{}.

\begin{quote}
\textbf{\zh{那个事故里有两个人重伤}} \emph{nàge shìgù lǐ yǒu liǎng ge rén zhòngshāng} — Two people were badly hurt in that accident.
\end{quote}

\begin{cousinweb}
\textbf{\zh{重伤}} (\emph{zhòngshāng}) is built from \zh{重} 'heavy' and \zh{伤} 'to hurt', both inherited from Middle Chinese, \emph{drjowngX} and \emph{syang} (Guangyun, in Baxter's transcription), which glosses the first 'thick' and the second 'to harm'.
\end{cousinweb}

\begin{grammarlens}[title={What you've built}]
One more word for saying how it was.
\end{grammarlens}

\subsection*{Your turn}
\begin{itemize}
\raggedright
  \item \emph{Say it:} \emph{zhòngshāng}
  \item \emph{Say it:} \emph{zhòngshāng}, once more
  \item \emph{Say it:} \emph{nàge shìgù lǐ yǒu liǎng ge rén zhòngshāng} — Two people were badly hurt in that accident.
  \item \emph{Recall:} say the Mandarin for to scald oneself, then the Mandarin for an open wound, then say \emph{zhòngshāng} again
\end{itemize}

\begin{tcolorbox}[breakable,colback=teal!4,colframe=teal!35!black,title={Before you move on}]
What does \zh{重伤} mean? (\textquotedblleft{}Serious injury; badly hurt\textquotedblright{}.) Say it once more.
\end{tcolorbox}
\section[niǔshāng]{niǔshāng — to sprain}

\label{lesson:ZH-C478-niushang}



\begin{quote}
Before the new one: say the Mandarin for an open wound, then the Mandarin for badly hurt.
\end{quote}

\subsection*{You'll want to know: niǔshāng}
\emph{niǔshāng} — \textquotedblleft{}to sprain\textquotedblright{}.

\begin{quote}
\emph{wǒ niǔshāng guo jiǎo} — I've sprained my ankle before.
\end{quote}

Its characters come later; for now it is a word for the ear.

\begin{cousinweb}
\emph{niǔshāng} is built from \emph{niǔ} 'to twist' and \zh{伤} 'to hurt', both inherited from Middle Chinese, \emph{nrjuwX} and \emph{syang} (Guangyun, in Baxter's transcription). The dictionary lists the first under two readings and glosses \emph{nrjuwX} as a hand turning.
\end{cousinweb}

\begin{grammarlens}[title={What you've built}]
One more word for saying how it was.
\end{grammarlens}

\subsection*{Your turn}
\begin{itemize}
\raggedright
  \item \emph{Say it:} \emph{niǔshāng}
  \item \emph{Say it:} \emph{niǔshāng}, once more
  \item \emph{Say it:} \emph{wǒ niǔshāng guo jiǎo} — I've sprained my ankle before.
  \item \emph{Recall:} say the Mandarin for an open wound, then the Mandarin for badly hurt, then say \emph{niǔshāng} again
\end{itemize}

\begin{tcolorbox}[breakable,colback=teal!4,colframe=teal!35!black,title={Before you move on}]
What does niǔshāng mean? (\textquotedblleft{}To sprain\textquotedblright{}.) Say it once more.
\end{tcolorbox}
\section[bāozā]{bāozā — to bandage; to dress a wound}

\label{lesson:ZH-C478-baoza}



\begin{quote}
Before the new one: say the Mandarin for badly hurt, then the Mandarin for to sprain.
\end{quote}

\subsection*{You'll want to know: bāozā}
\emph{bāozā} — \textquotedblleft{}to bandage; to dress a wound\textquotedblright{}.

\begin{quote}
\emph{hùshi gěi wǒ bāozā le shāngkǒu} — The nurse dressed my wound.
\end{quote}

Its characters come later; for now it is a word for the ear.

\begin{cousinweb}
\emph{bāozā} is built from \zh{包} 'to wrap' and \emph{zā} 'to bind', both inherited from Middle Chinese, \emph{paew} and \emph{tsreat} (Guangyun, in Baxter's transcription), which glosses the first 'to wrap' and the second, in its fuller spelling, as binding the grip of a bow. The regular sound change would give \emph{zhá}; why it is said \emph{zā} is not settled.
\end{cousinweb}

\begin{grammarlens}[title={What you've built}]
One more word for saying how it was.
\end{grammarlens}

\subsection*{Your turn}
\begin{itemize}
\raggedright
  \item \emph{Say it:} \emph{bāozā}
  \item \emph{Say it:} \emph{bāozā}, once more
  \item \emph{Say it:} \emph{hùshi gěi wǒ bāozā le shāngkǒu} — The nurse dressed my wound.
  \item \emph{Recall:} say the Mandarin for badly hurt, then the Mandarin for to sprain, then say \emph{bāozā} again
\end{itemize}

\begin{tcolorbox}[breakable,colback=teal!4,colframe=teal!35!black,title={Before you move on}]
What does bāozā mean? (\textquotedblleft{}To bandage; to dress a wound\textquotedblright{}.) Say it once more.
\end{tcolorbox}
